\documentclass{article} %210 mm × 297 mm
\usepackage[utf8]{inputenc}
\usepackage[a4paper,left=1.5cm, right=1.5cm, top=2cm, bottom=2cm]{geometry}
\usepackage{graphicx}
%\usepackage{blindtext}
\usepackage{multirow}
\usepackage{mhchem}
\usepackage{url}
\usepackage{subfigure}
\usepackage{amsfonts,latexsym} % para tener disponibilidad de diversos simbolos
\usepackage{enumerate}
%\usepackage{booktabs}
\usepackage{float}
%\usepackage{threeparttable}
\usepackage{array,colortbl}
%\usepackage{ifpdf}
%\usepackage{rotating}
\usepackage{cite}
\usepackage{stfloats}
\usepackage{url}
\usepackage{listings}
\usepackage{fancyhdr}
\usepackage{biblatex}
\usepackage{color}
\addbibresource{sources.bib}
%\usepackage{hyperref}
%\usepackage{wrapfig}
\usepackage{array}
%\usepackage{multirow}
%\usepackage{adjustbox}
%\usepackage{nccmath}
\usepackage[dvipsnames]{xcolor}
\usepackage{tcolorbox}
\newtcolorbox{mybox}{colback=white,
colframe=black}
\newcommand{\titulo}{Abgabe 4; Diskrete Mathematik}
\newcommand{\nombre}{Liliana Sanfilippo}
\newcommand{\FCE}{\textbf{WiSe 2024/2025}}

 

\usepackage{fancyhdr}
\pagestyle{fancy}
\fancyhead{}
\fancyfoot{}
\fancyhead[LO]{\small\nombre}
\fancyhead[RE]{\small\titulo}
\fancyhead[LO]{\small\FCE}
\fancyfoot[RO,LE] {\thepage}
\usepackage[onehalfspacing]{setspace}

\setcounter{page}{1} %Número de la página, la editorial lo cambiará
\begin{document}
\begin{tabular}{p{12cm}}
{{\Large\bf\titulo}}\\
\vspace{0.2cm}\\
 {\large\nombre}\\
 \vspace{0.2cm}\\   
\end{tabular}
\section*{Aufgaben zur Abgabe}
\subsection*{Aufgabe 1 (4 Punkte)}
Wie viele verschiedene Wörter kann man durch Permutation der Buchstaben aus den
Buchstaben \( A, E, M, O, U, Y \) bilden? In wie vielen dieser Wörter taucht weder \texttt{ME} noch \texttt{YOU} auf?
\begin{mybox}
 \(| A, E, M, O, U, Y | = 6\) und $6! = 720$, also gibt es 720 verschiedene Wörter (das bezieht sich aber nicht auf Wörter mit semantischem Sinn!). \\ 
 Um herauszufinden in wie vielen dieser Wörter weder \texttt{ME} noch \texttt{YOU} auftaucht, kann man diese Kombinationen als Buchstaben betrachten.  Zumindest wenn man es als Permutationen ohne Wiederholung betrachtet. \\
  \(| A, ME, O, U, Y | = 5\) und $5! = 120$. \\
   \(| A, E, M, YOU | = 4\) und $4! = 24$. \\
      \(| A, ME YOU | = 3\) und $3! = 6$. \\
      Mit der Inklusions-Exklusionsregel: 
       \[
\text{Wörter ohne ME und YOU} = \text{Gesamtanzahl} -( \text{alle Wörter mit ME} + \text{alle Wörter mit YOU} - |A \cap B|)
\] also
      \[
\text{Wörter ohne ME und YOU} = \text{Gesamtanzahl} - \text{alle Wörter mit ME} - \text{alle Wörter mit YOU} + |A \cap B|
\]
      \[
      720-120-24+6=582
      \]
\end{mybox}
\subsection*{Aufgabe 2 (4 Punkte)}
Eine Permutation \( \sigma \in S_n \) heißt zyklisch, falls sie durch einen Zyklus der Länge \( n \) dargestellt
werden kann. Zeigen Sie, dass es \( (n - 1)! \) zyklische Permutationen gibt.
\begin{mybox}
Man kann jede der \( (n - 1)! \) zyklischen Permutationen darstellen, indem man irgendeine der n Positionen als Start des Zyklus wählt und die restlichen Elemente (n-1) in einer bestimmten Reihenfolge ordnet. Also gibt es \( (n - 1)! \) zyklische Permutationen.
\end{mybox}
\subsection*{Aufgabe 3 (4 Punkte)}
Finden Sie für jede der folgenden Werte \( (v, k, r) \) ein Design mit diesen Parametern oder
erklären Sie, warum ein solches nicht existiert:
\begin{itemize}
    \item[(a)] \( (6, 3, 2) \)
    \item[(b)] \( (5, 2, 1) \)
    \item[(c)] \( (7, 4, 4) \)
    \item[(d)] \( (9, 6, 4) \)
\end{itemize}
\begin{mybox}
$b = \frac{r \cdot v}{k}$
\begin{itemize}
    \item[(a)] \((v, k, r) = (6, 3, 2) \) \\
    Gilt $k | v\cdot r$ ? \\
    Ja, $4 | 6\cdot 2$ also $4 | 12 $ gilt.  \\
    $b = \frac{2 \cdot 6}{3} = 4$ \\
    Gilt $b \leq \binom{v}{k}$? \\
    $
    \binom{v}{k} = \binom{6}{3} = 20
    $ also ja, $4 \leq 20$ ist richtig. \\
    Also gibt es ein solches Design. 
    \item[(b)] \((v, k, r) =  (5, 2, 1) \)
    $b = \frac{1 \cdot 5}{2} = $ \\
    Gilt $2 | 5\cdot 1$ ? \\
    Nein, 2 ist kein Teiler von 5. Daher kann es kein solches Design geben. 
    
    \item[(c)] \((v, k, r) =  (7, 4, 4) \) \\
    Gilt $4 | 7\cdot 4$ ? \\
    Ja, 4 ist ein Teiler von 28.  \\
    $b = \frac{4 \cdot 7}{4} = 7$ \\
    Gilt $b \leq \binom{v}{k}$? \\ 
    Ja, denn $7 \leq 35$ ist richtig. 
    Also gibt es ein solches Design. 
    
    \item[(d)] \((v, k, r) =  (9, 6, 4) \) \\
    Gilt $6 | 9\cdot 4$ ? \\ 
    Ja, 6 ist ein Teiler von 36. \\
    $b = \frac{4 \cdot 9}{6} = 6$ \\
     Gilt $b \leq \binom{v}{k}$? \\
     Ja, denn $6 \leq 84$
     Also gibt es ein solches Design. 
     
\end{itemize}
\end{mybox}
\subsection*{Aufgabe 4 (4 Punkte)}
Ein 2-Design mit Parametern \( (v, k, \lambda) = (v, 3, 1) \) wird \emph{Steinersches Tripelsystem} (STS) genannt.
Stellen Sie fest, für welche \( v \) im Bereich \( 3 \leq v \leq 12 \) ein STS mit \( v \) Varietäten
existieren kann, und finden Sie ein Beispiel in diesen Fällen.
\begin{mybox}
Absolut gar keine Ahnung, ich verstehe STS anscheinend nicht. 
\end{mybox}



\end{document}